% ---------------------------------------------------------------------------
% Chapter -- Theoretical Background and Related Work
% Text rewritten in plain English (1 Nov 2026). Statements about other papers
% are NOT all verified against the papers; every open check is marked \TODO.
% ---------------------------------------------------------------------------

\chapter{Theoretical Background and Related Work}
\label{ch:background}

This chapter collects the ideas on which the rest of the thesis is built.
Section~\ref{sec:bg-constraints} explains what makes a system nonholonomic.
Section~\ref{sec:bg-brackets} shows how Lie brackets decide controllability, and
how ``deep'' a system is. Section~\ref{sec:bg-brockett} explains why no smooth
time-invariant feedback can stabilize such a system.
Section~\ref{sec:bg-classical} describes the classical tools for vehicles with
trailers and why they do not cover off-hooked trailers.
Section~\ref{sec:bg-oscillating} introduces the oscillating controllers used in
this thesis, and Section~\ref{sec:bg-state} places the thesis in the literature.

% ===========================================================================
\section{Nonholonomic constraints}
\label{sec:bg-constraints}

A wheel that rolls without slipping cannot move sideways. Take a planar body
with position $(x,y)$ and heading $\theta$, carried by a wheel axle with its
midpoint at $(x,y)$. The condition of no sideways motion is
\begin{equation}
  \dot x\sin\theta-\dot y\cos\theta=0 .
  \label{eq:bg-rolling}
\end{equation}
This is a linear constraint on the velocity $\dot q$ of the configuration
$q=(x,y,\theta)$. Constraints of this kind are written in Pfaffian form
$A(q)\dot q=0$ \cite{deluca1995}.

The word ``constraint'' can mean two different things. If
$A(q)\dot q=0$ can be integrated to a relation $h(q)=0$ between the
coordinates, the constraint is \emph{holonomic}. It removes configurations, and
the system lives on a surface of lower dimension. If it cannot be integrated,
the constraint is \emph{nonholonomic}. It removes only \emph{directions of
instantaneous motion}, and every configuration may still be reachable. The
wheel is the standard example. At any instant it cannot move sideways, yet a
car can be parked in a gap next to it, because forward, turning and backward
motions together produce a net sideways displacement. The constraint limits how
the system moves, not where it can go.

The admissible velocities are the vectors in the kernel of $A(q)$. A basis
$f_1(q),\dots,f_m(q)$ of this kernel turns the constraint into a
\emph{driftless control-affine system}
\begin{equation}
  \dot q=\sum_{i=1}^{m}u_if_i(q),\qquad q\in\R^n,\ u\in\R^m,\ m<n,
  \label{eq:bg-driftless}
\end{equation}
where the inputs $u_i$ are velocities. For~\eqref{eq:bg-rolling} one can take
\begin{equation}
  f_1(q)=\begin{pmatrix}\cos\theta\\ \sin\theta\\ 0\end{pmatrix},\qquad
  f_2(q)=\begin{pmatrix}0\\ 0\\ 1\end{pmatrix},
  \label{eq:bg-unicycle}
\end{equation}
with $u_1$ the forward velocity and $u_2$ the angular velocity. This is the
\emph{unicycle}, the kinematic model of a differential-drive robot such as the
TurtleBot3 used in this thesis. Each trailer adds one angle and one rolling
constraint. A car with $N$ trailers therefore has $n=N+3$ states and still only
$m=2$ inputs. The model for $N=1,2$ off-hooked trailers is derived in
Chapter~\ref{ch:model}. ``Driftless'' means that $u=0$ gives $\dot q=0$:
the vehicle stops when the wheels stop.

Whether a constraint is holonomic is decided by the Frobenius theorem. The
constraints are integrable exactly when the distribution spanned by
$f_1,\dots,f_m$ is involutive, that is, closed under the Lie bracket introduced
next \cite{nijmeijer1990,deluca1995}. For the unicycle it is not, as
\cref{eq:bg-unicycle-bracket} below shows.

% ===========================================================================
\section{Lie brackets and controllability}
\label{sec:bg-brackets}

\begin{definition}[Lie bracket]
For smooth vector fields $f,g$ on $\R^n$ the \emph{Lie bracket} is the vector
field
\begin{equation}
  \bracket{f}{g}(q)=\frac{\partial g}{\partial q}(q)\,f(q)
                   -\frac{\partial f}{\partial q}(q)\,g(q).
  \label{eq:bg-bracket}
\end{equation}
\end{definition}

The bracket has a direct meaning for motion. Let $\Phi^f_t$ be the flow of $f$
for time $t$. The system moves along $f$ for time $t$, then along $g$, then
along $-f$, and finally along $-g$. It does not return to $q$, but to
\begin{equation}
  \Phi^{-g}_t\circ\Phi^{-f}_t\circ\Phi^{g}_t\circ\Phi^{f}_t(q)
  =q+t^2\,\bracket{f}{g}(q)+O(t^3).
  \label{eq:bg-commutator}
\end{equation}
A closed loop in the inputs therefore produces a net motion in the direction
of the bracket. The size of this motion is $t^2$ instead of $t$. For the unicycle~\eqref{eq:bg-unicycle}
\begin{equation}
  \bracket{f_1}{f_2}(q)=\begin{pmatrix}\sin\theta\\ -\cos\theta\\ 0\end{pmatrix},
  \label{eq:bg-unicycle-bracket}
\end{equation}
which is exactly the sideways direction forbidden by~\eqref{eq:bg-rolling}.
The manoeuvre ``drive, turn, drive back, turn back'' is the parking manoeuvre of
Section~\ref{sec:bg-constraints}. Equation~\eqref{eq:bg-commutator} says how much it
achieves: driving and turning for a time $t$ each moves the robot sideways
by about $t^2$. The bracket is not in the span of $f_1,f_2$. The distribution
is therefore not involutive, and the constraint~\eqref{eq:bg-rolling} is
nonholonomic.

Brackets of brackets give further directions, which are reached by longer
manoeuvres. Following \cite{bellaiche2005}, let $\mathcal L^1$ be the span of
$f_1,\dots,f_m$, and let $\mathcal L^s$ be the span of all brackets of these
fields of length at most $s$. A bracket such as $\bracket{f_1}{\bracket{f_1}{f_2}}$
has length~$3$. At a point $p$, these spaces form an increasing sequence of
subspaces $\mathcal L^1(p)\subseteq\mathcal L^2(p)\subseteq\dots\subseteq\R^n$.

\begin{definition}[Chow's condition, LARC]
System~\eqref{eq:bg-driftless} satisfies \emph{Chow's condition}, also called
the \emph{Lie algebra rank condition} (LARC), if the fields $f_i$ and their
iterated brackets span $\R^n$ at every point.
\end{definition}

\begin{theorem}[Chow, Rashevskii {\cite[Thm.~1.3]{bellaiche2005}}]
If the state space is connected and Chow's condition holds, any two points can
be joined by a trajectory of~\eqref{eq:bg-driftless}.
\end{theorem}

The theorem gives controllability in the sense of reachability. It says nothing
about how to steer or how long it takes. When Chow's condition holds, there is
at every point $p$ a smallest integer $r(p)$ with $\mathcal L^{r(p)}(p)=\R^n$.
This is the \emph{degree of nonholonomy} at $p$. The dimensions
\begin{equation}
  \bigl(n_1(p),n_2(p),\dots,n_r(p)\bigr),\qquad n_s(p)=\dim\mathcal L^s(p),
\end{equation}
form the \emph{growth vector} \cite{bellaiche2005}. A point is \emph{regular}
if these dimensions are constant in a neighbourhood of it, and \emph{singular}
otherwise \cite{bellaiche2005}.

\paragraph{Physical reading of the growth vector.}
The first $n_1$ directions can be driven directly by the inputs. The next
$n_2-n_1$ directions need a manoeuvre of the type~\eqref{eq:bg-commutator},
whose effect after time $t$ is of order $t^2$. The next $n_3-n_2$ directions
need manoeuvres built from brackets of length three, with an effect of order
$t^3$, and so on. The degree of nonholonomy is thus the length of the
longest manoeuvre needed. Directions of high bracket length are
``expensive'': a small motion in them needs a comparatively large motion
of the inputs. This is the origin of the amplitude problem quantified in
Chapter~\ref{ch:feasibility}.

\paragraph{Vehicles with trailers.}
Consider the standard car with trailers, hitched at the axle midpoints.
Bella\"iche, Jean and Risler show that the car with one trailer has growth
vector $(2,3,4)$ everywhere, so $r=3$. The car with two trailers has $r=4$ at
regular points and $r=5$ at the singular points where two consecutive bodies are
perpendicular \cite{bellaiche2005}. For a car with two \emph{off-hooked}
trailers and unit lengths, Vendittelli, Laumond and Oriolo find $r=3$ at
regular points and $r=4$ at singular points. The singular points are those
where the two hitch angles are equal \cite{vendittelli1998}. The straight
configuration, which is the goal of this thesis, has both hitch angles equal to
zero. It is therefore singular in this sense. The value $r=4$ at the origin,
obtained in~\cite{zuyev2024} and in Chapter~\ref{ch:model} for arbitrary lengths,
is consistent with this. The goal therefore lies exactly where the system is
hardest to control.
\TODO[Check in bellaiche2005 and vendittelli1998 the values of r at regular and singular points quoted above; only the titles of these papers were found]

% ===========================================================================
\section{Brockett's condition}
\label{sec:bg-brockett}

Controllability does not imply that a feedback law $u=k(q)$ can stabilize an
equilibrium. Brockett showed this with a degree-theoretic argument
\cite{brockett1983}.

\begin{theorem}[Brockett {\cite[Thm.~1]{brockett1983}}]
\label{thm:brockett}
Let $\dot q=f(q,u)$ with $f(q_0,0)=0$ and $f$ continuously differentiable near
$(q_0,0)$. If there is a continuously differentiable feedback $u=k(q)$ that
makes $q_0$ asymptotically stable, then
\begin{enumerate}
  \item the linearization has no uncontrollable modes with eigenvalues of
        positive real part,
  \item there is a neighbourhood $N$ of $q_0$ such that every $q\in N$ can be
        steered to $q_0$ as $t\to\infty$, and
  \item the map $(q,u)\mapsto f(q,u)$ is onto an open set containing $0$.
\end{enumerate}
\end{theorem}

The third condition is the restrictive one. It asks that every small velocity
$\xi$ can be produced by some state near $q_0$ and some small input. For a
driftless system~\eqref{eq:bg-driftless} whose fields $f_1,\dots,f_m$ are
linearly independent at $q_0$, this fails as soon as $m<n$. Take coordinates in
which $f_i(q_0)$ is the $i$-th unit vector. Then the first $m$ components of
$\sum_iu_if_i(q)$ vanish for $q$ near $q_0$ only if $u=0$. So no small velocity
of the form $(0,\dots,0,\xi)$ with $\xi\neq0$ is produced. Brockett concludes
that such a system can be stabilized by continuously differentiable feedback if
and only if $m=n$ \cite{brockett1983}. Every nonholonomic vehicle model, with
$m<n$, is therefore excluded. Brockett's own example is the system
$\dot x=u,\ \dot y=v,\ \dot z=xv-yu$, now often called the Brockett
integrator \cite{brockett1983}. For the vehicle of this thesis, the violation of
the third condition is shown directly in Chapter~\ref{ch:model}. A related
result of Artstein, published in the same year, shows that systems which
cannot be stabilized by continuous feedback may still be stabilized with
relaxed controls \cite{artstein1983}.

Theorem~\ref{thm:brockett} leaves three ways out, and all are used in the
literature:
\begin{itemize}
  \item \emph{Time-varying feedback} $u=k(t,q)$. Coron proved that a
        controllable driftless system can be asymptotically stabilized by a
        smooth feedback that is periodic in time \cite{coron1992}. The
        controllers of this thesis belong to this class.
        \TODO[Check the exact conditions of Coron's theorem (smoothness, global or local) in coron1992]
  \item \emph{Discontinuous feedback}, which avoids the continuity assumption.
  \item \emph{A weaker goal}: convergence of only part of the state, or to a
        neighbourhood of the equilibrium (practical stabilization).
\end{itemize}

% ===========================================================================
\section{Classical approaches for vehicles with trailers}
\label{sec:bg-classical}

\paragraph{Chained form.}
Many two-input nonholonomic systems can be transformed, by a change of
coordinates and inputs, into the \emph{chained form}
\begin{equation}
  \dot z_1=v_1,\quad \dot z_2=v_2,\quad \dot z_3=z_2v_1,\quad\dots,\quad
  \dot z_n=z_{n-1}v_1 .
  \label{eq:bg-chained}
\end{equation}
For this form, steering with sinusoidal inputs \cite{murray1993} and
exponential stabilization by time-varying homogeneous feedback
\cite{mccloskey1997} are available. This makes the transformation attractive.
Sørdalen showed that the kinematics of a car with $n$ trailers, each hitched at
the axle midpoint of the preceding body, can be converted into chained form
\cite{sordalen1993}. The conversion is local. As noted in
\cite{michalek2014tcst}, the transformation of \cite{sordalen1993} is valid
only if the orientation of the last trailer and all joint angles stay in
$(-\pi/2,\pi/2)$. According to the experience reported in
\cite{mccloskey1994}, such strongly nonlinear transformations are
ill-conditioned, and the resulting controller can be very sensitive to small
changes of the configuration \cite{michalek2014tcst}. M'Closkey and Murray
implemented time-periodic exponential stabilizers on a mobile robot towing a
trailer, and compared several control laws experimentally
\cite{mccloskey1994}.
\TODO[Check in mccloskey1994 how the trailer was hitched (on or off the axle) and the robot used]

\paragraph{Flatness.}
A system is differentially flat if all states and inputs can be written as
functions of a set of outputs and finitely many of their derivatives. Motion
planning then reduces to choosing a smooth curve for these outputs. The
standard car with trailers is flat \cite{rouchon1993}, with the position of
the last trailer as flat output \cite{michalek2014tcst}.
\TODO[Check whether the car with ONE off-hooked trailer is flat (Rouchon et al. 1993 may treat it); if so, say that flatness-based planning exists for r = 3 but not for two off-hooked trailers]

\paragraph{Why both fail for off-hooked trailers.}
With the hitch off the axle, the car with two trailers cannot be put into
chained form. It also does not satisfy the necessary and sufficient flatness
condition for two-input driftless systems \cite{vendittelli1998}. This is also
the starting point of \cite{zuyev2024}. Michałek and co-workers describe
$N$-trailers with off-axle hitches as \emph{differentially non-flat}
\cite{michalek2015jint}. Altafini interprets the additional complexity as
virtual steering wheels at the off-axle joints. Their steering angles are
nonlinear feedbacks of the configuration. This is used to embed the general
$n$-trailer into a multi-steering $n$-trailer, for which a multi-chained form is
available \cite{altafini2001}. Lizarraga, Morin and Samson use a chained-form
\emph{approximation} of the general $N$-trailer for exponential stabilization
\cite{lizarraga2001}.
\TODO[Read lizarraga2001: is the result local, practical or exact, and is it validated by simulation or experiment?]

Chapter~\ref{ch:nilpotent} shows the connection from the other side. In the
limit $\dzero\to0$, the nilpotent approximation of the off-hooked system becomes
the chained form. The off-axle distance enters through a single additional term.

% ===========================================================================
\section{Oscillating controls from gradient-flow approximations}
\label{sec:bg-oscillating}

The idea behind the controllers of this thesis is~\eqref{eq:bg-commutator} used
in reverse. A Lyapunov function $V$ should decrease. The ideal velocity for this
is the negative gradient $-\nabla V$. In general it points in directions that
are reachable only through brackets. A periodic input of small period
$\varepsilon$ approximates this gradient flow, if its oscillating parts are
chosen so that their net effect over one period equals a prescribed combination
of the fields and their brackets. The open-loop construction is turned into a
feedback by \emph{sampling}. At the beginning of each period the state is
measured, the coefficients are computed from it, and the input of that period
is applied.

The method was developed in a sequence of papers by Zuyev, Grushkovskaya and
co-authors:
\begin{itemize}
  \item Zuyev \cite{zuyev2015} constructs trigonometric open-loop controls that
        approximate the gradient flow of a Lyapunov function. Under the
        sampling strategy, exponential stability is proved for driftless
        systems whose fields and first-order brackets span the tangent space.
        The example is the Brockett integrator.
  \item Zuyev, Grushkovskaya and Benner \cite{zuyev2016} extend this to
        systems that need brackets of second order.
  \item Grushkovskaya and Zuyev \cite{grushkovskaya2023} give a unified
        framework for stabilization, motion planning and obstacle avoidance.
        It is based on gradient-flow approximation under bracket-generating
        conditions.
  \item Zuyev and Grushkovskaya \cite{zuyev2024} apply the scheme to a car
        with two off-hooked trailers ($r=4$). They derive privileged coordinates
        and the nilpotent approximation explicitly in terms of the four
        lengths, design an oscillating controller and prove exponential
        convergence. The convergence result is limited to the nilpotent
        approximation, and the original system is left to future work. Their
        simulations use the approximation.
  \item Zuyev and Grushkovskaya \cite{zuyev2025classical} give sufficient
        conditions under which classical solutions of the closed loop, not
        only the sampled ones, converge.
  \item Zuyev, Grushkovskaya and Eisner \cite{zuyev2026implementation}
        implement controllers of this family on a TurtleBot3 Burger, a unicycle
        ($r=2$), with OptiTrack motion capture. They report chattering caused by
        the oscillating inputs under the velocity limits of the robot. They study
        the admissibility of the gradient flow, and show that the reference
        position can be stabilized with practically acceptable parameters.
\end{itemize}
Murray and Sastry's steering with sinusoids \cite{murray1993} uses the same
mechanism for open-loop motion planning in chained form. The difference here is
that the oscillations are generated in feedback, from the current state, and
that the target is the full configuration.

The controller is described in detail in Chapter~\ref{ch:control}. Two of its
properties matter for what follows. First, the size of an oscillation that
produces a drift $a$ in a direction of bracket length $w$ grows like $|a|^{1/w}$.
The deep directions are therefore expensive, as anticipated in
\cref{sec:bg-brackets}. Second, the convergence proof holds only for a
sufficiently small period $\varepsilon$, which depends on the size of the set of
initial states, and only for the nilpotent approximation. How the controller
behaves on the exact kinematics and within the velocity limits of a real robot
is not covered by the theory.

% ===========================================================================
\section{State of the art}
\label{sec:bg-state}

\paragraph{Terminology.}
The name ``off-hooked trailers'' is used by Zuyev and Grushkovskaya
\cite{zuyev2024} and by Vendittelli et al.\ \cite{vendittelli1998}. Most of the
literature uses \emph{off-axle hitching}. It distinguishes standard N-trailers
(all hitches on the axle), non-standard N-trailers (all hitches off the axle)
and general N-trailers (both kinds)
\cite{michalek2014pf,michalek2014tcst}. The term \emph{kingpin hitch} is also
used \cite{altafini2001}. In this terminology the vehicle of this thesis is a
non-standard N-trailer with $N=1$ or $N=2$.

\paragraph{Tasks.}
The control tasks must be separated, because their difficulty is very
different. In \emph{path following}, the vehicle has to move along a given
geometric path. In \emph{trajectory tracking}, it has to follow a given
time-dependent reference. Both can be solved with time-invariant feedback,
because the reference keeps moving. \emph{Set-point control} (point
stabilization) asks the vehicle to come to rest at a given configuration. This
is the task affected by Theorem~\ref{thm:brockett}. Within set-point control,
steering only a guidance point of the vehicle, for example the last trailer, is
weaker than stabilizing the full configuration.

\paragraph{Related work.}
Michałek's group in Poznań has the most extensive work on N-trailers with
off-axle hitching, including hardware experiments with a three-trailer
laboratory vehicle. Their path-following controller for N-trailers with
off-axle hitching \cite{michalek2014pf} was validated experimentally on a
three-trailer vehicle. So was their cascaded Vector-Field-Orientation (VFO)
controller for set-point and trajectory tracking of non-standard N-trailers
\cite{michalek2015jint}. The latter takes the input limits of the robot into
account by a velocity-scaling procedure. A cascaded VFO set-point controller
for the on-axle case was also verified by simulations and experiments
\cite{michalek2014tcst}. An earlier paper applies the VFO method to set-point
control of the N-trailer with off-axle hitching \cite{michalek2012}.
\TODO[Michalek 2012 (IJC 85(5):502--521): only the title is verified. Check in the full text what is guaranteed (whole configuration or guidance point, exact or neighbourhood) and whether it is simulation only]
The same authors write that, for non-standard and general N-trailers, a
practically acceptable solution of set-point stabilization of an
\emph{arbitrary configuration of the whole vehicle} ``seems to be still open''
\cite{michalek2015jint}.

\Cref{tab:bg-state} summarizes the approaches by task, system and validation.

\begin{table}[ht]
  \centering
  \small
  \caption{Approaches to the control of vehicles with trailers. ``Full state''
    means stabilization of the whole configuration. Entries marked $^\ast$ are
    taken from titles or abstracts only and still have to be checked.}
  \label{tab:bg-state}
  \begin{tabular}{@{}>{\raggedright\arraybackslash}p{3.6cm}>{\raggedright\arraybackslash}p{3.1cm}>{\raggedright\arraybackslash}p{2.7cm}>{\raggedright\arraybackslash}p{3.2cm}@{}}
    \toprule
    approach & task & hitching & validation \\
    \midrule
    chained form \cite{sordalen1993,mccloskey1997} & set-point, full state
      & on-axle & experiments with one trailer$^\ast$ \cite{mccloskey1994} \\
    chained-form approximation \cite{lizarraga2001} & set-point (exponential)
      & general N-trailer & $^\ast$ \\
    VFO path following \cite{michalek2014pf} & path following & off-axle
      & experiments, 3 trailers \\
    cascaded VFO \cite{michalek2014tcst} & set-point & on-axle
      & simulations, experiments, 3 trailers \\
    cascaded VFO \cite{michalek2015jint} & set-point, tracking & off-axle
      & experiments, 3 trailers \\
    VFO set-point \cite{michalek2012} & set-point & off-axle & $^\ast$ \\
    nilpotent approx.\ \cite{vendittelli1998} & steering (open loop)
      & off-axle, 2 trailers & simulations \cite{zuyev2024} \\
    oscillating control \cite{zuyev2024} & set-point, full state
      & off-axle, 2 trailers & simulations on the approximation \\
    oscillating control \cite{zuyev2026implementation} & set-point
      & unicycle (no trailer) & experiments \\
    \textbf{this thesis} & set-point, full state & off-axle, 1 and 2 trailers
      & exact kinematics; Gazebo and experiments planned \\
    \bottomrule
  \end{tabular}
\end{table}

\paragraph{Position of this thesis.}
The thesis takes the oscillating controller of \cite{zuyev2024} and moves it
towards a real vehicle. The model and the nilpotent approximation are derived
for one off-hooked trailer ($r=3$), which \cite{zuyev2024} does not treat. The
controller is studied on the exact kinematics instead of the approximation, and
within the velocity limits of the TurtleBot3. Finally, it is validated in
Gazebo and in experiments.
\TODO[Gazebo and the experiments are not done yet; adjust this sentence and the table row when they are]
Experiments with time-varying stabilizers for a robot with a trailer exist
\cite{mccloskey1994}, and experiments with oscillating controllers of this
family exist for the unicycle \cite{zuyev2026implementation}. In the literature found for this thesis, there is no experimental validation of
the oscillating gradient-flow controller for off-hooked trailers, or for a
degree of nonholonomy above two.
\TODO[Confirm the novelty statement with the supervisor; the second half depends on what mccloskey1994 and lizarraga2001 actually do]
